\section{Описание}

В работе реализован алгоритм Ахо--Корасик для поиска множества образцов в тексте. Алгоритм сначала
строит trie по всем образцам, затем дополняет его связями неудачи и выходными ссылками. После этого
текст обрабатывается за один проход: каждое очередное число переводит автомат в новое состояние, а
все образцы, заканчивающиеся в этом состоянии, сразу выводятся как найденные \cite{AhoCorasick, Kormen3}.

Особенность варианта состоит в большом алфавите: символами являются не буквы, а числа типа
\texttt{uint32\_t}. Поэтому переходы из вершины нельзя хранить в массиве размера $2^{32}$.
В реализации используется разреженное хранение переходов:

\begin{center}
\texttt{std::unordered\_map<std::uint32\_t, std::size\_t>}
\end{center}

Такой подход хранит только реально существующие рёбра trie и соответствует требуемому алфавиту.
Записи вида \texttt{0002}, \texttt{02} и \texttt{2} при чтении превращаются в одно и то же число.

\subsection{Структура автомата}

Вершины автомата хранятся в массиве \texttt{nodes\_}. Корень имеет индекс \(0\). Каждая вершина содержит:
\begin{itemize}
\item \texttt{next} -- переходы по числам алфавита;
\item \texttt{failure\_link} -- связь неудачи, то есть переход к наибольшему собственному суффиксу текущего состояния;
\item \texttt{output\_link} -- ссылка на ближайшую вершину по цепочке связей неудачи, в которой заканчивается образец;
\item \texttt{patterns} -- список образцов, которые заканчиваются непосредственно в данной вершине.
\end{itemize}

Для найденного образца хранится его номер во входе и длина. Длина нужна не для работы автомата, а для
восстановления позиции начала вхождения в тексте.

\subsection{Добавление образцов}

Метод \texttt{AddPattern} добавляет один образец в trie. Обход начинается из корня. Для каждого числа
образца проверяется наличие перехода из текущей вершины. Если переход отсутствует, создаётся новая
вершина; если переход уже существует, обход продолжается по нему. После обработки последнего числа в
текущей вершине сохраняется информация о завершившемся образце: его номер и длина.

Также при добавлении поддерживается значение \texttt{max\_pattern\_length\_}. Оно равно максимальной
длине среди всех образцов и далее используется для хранения только нужного числа последних позиций текста.

\subsection{Построение связей}

Метод \texttt{BuildLinks} строит связи неудачи обходом в ширину. Для детей корня связь неудачи ведёт в
корень. Для каждой другой вершины рассматривается ребро из родителя по некоторому числу \texttt{value}.
Алгоритм берёт связь неудачи родителя и поднимается по цепочке связей неудачи, пока не найдёт вершину,
из которой существует переход по \texttt{value}. Если такой переход найден, связь неудачи текущего
ребёнка указывает на соответствующую вершину; иначе она указывает на корень.

После связи неудачи для вершины строится \texttt{output\_link}. Если вершина по связи неудачи сама
содержит завершение образца, выходная ссылка ведёт в неё. Иначе используется выходная ссылка этой
вершины. Это позволяет быстро находить все образцы, являющиеся суффиксами текущей прочитанной
последовательности.

\subsection{Переходы при поиске}

Метод \texttt{NextState} выполняет переход автомата по очередному числу текста. Если из текущей вершины
есть переход по этому числу, автомат сразу переходит по нему. Если перехода нет, выполняется движение по
\texttt{failure\_link}, пока переход не будет найден или пока автомат не вернётся в корень. Если перехода
нет даже из корня, состояние остаётся равным нулю.

Метод \texttt{ForEachMatch} проверяет все образцы, которые заканчиваются в текущем состоянии. Сначала
рассматриваются образцы самой вершины, затем выполняются переходы по \texttt{output\_link}. Для каждого
найденного образца вызывается переданная callback-функция, которая в \texttt{main} печатает ответ.

\subsection{Чтение входа и восстановление позиции}

Функция \texttt{ParseNumbers} разбирает одну строку на числа и для каждого числа вызывает переданный
обработчик. Такой формат позволяет использовать один и тот же парсер и для строк с образцами, и для
строк текста.

Во время обработки текста автомат находит момент окончания образца. В ответе же требуется вывести
позицию начала. Для этого используется очередь \texttt{recent\_positions}, где хранятся позиции последних
прочитанных чисел. Достаточно хранить не больше \texttt{max\_pattern\_length} позиций: более старое число
не может быть началом только что найденного образца.

Если найден образец длины \(k\), то его начало находится на \(k\)-й позиции с конца очереди:

\begin{center}
\texttt{recent\_positions.size() - pattern.length}
\end{center}

Полученная позиция выводится вместе с номером найденного образца.

Асимптотика реализации:
\begin{itemize}
\item построение trie -- \(O(P)\), где \(P\) -- суммарная длина всех образцов;
\item построение связей -- \(O(P)\) в среднем при использовании \texttt{unordered\_map};
\item поиск -- \(O(T + Ans)\), где \(T\) -- число чисел в тексте, \(Ans\) -- число найденных вхождений;
\item память -- \(O(P)\) на автомат и \(O(L)\) на очередь позиций, где \(L\) -- максимальная длина образца.
\end{itemize}

\section{Исходный код}

Ниже приведены основные фрагменты реализации.

\lstinputlisting[
language=C++,
firstline=10,
lastline=27,
caption={Структуры позиций, образцов и вершины автомата},
label={lst:node}
]{../main.cpp}

\lstinputlisting[
language=C++,
firstline=30,
lastline=47,
caption={Добавление образца в trie},
label={lst:add-pattern}
]{../main.cpp}

\lstinputlisting[
language=C++,
firstline=49,
lastline=86,
caption={Построение связей неудачи и выходных ссылок},
label={lst:build-links}
]{../main.cpp}

\lstinputlisting[
language=C++,
firstline=92,
lastline=113,
caption={Переход автомата и перечисление найденных образцов},
label={lst:search-methods}
]{../main.cpp}

\lstinputlisting[
language=C++,
firstline=120,
lastline=196,
caption={Разбор ввода и потоковая обработка текста},
label={lst:main}
]{../main.cpp}
\pagebreak

\section{Консоль}
\begin{alltt}
$ g++ -std=c++14 -O2 -Wall -Wextra -pedantic main.cpp -o lab4.out
$ ./lab4.out
1 0002 001 2
1 2 1
2 02 2 2

01 2 1 2   1 2 1 3
2 2 2 2 02
1, 1, 2
1, 1, 1
1, 3, 2
1, 3, 1
1, 5, 2
2, 1, 3
2, 2, 3
\end{alltt}

Порядок вывода может отличаться от порядка в примере из условия, поскольку в условии явно указано,
что порядок следования найденных вхождений несущественен.
\pagebreak
